\chapter{Exemplos de aplicação direta}
\label{cap:exemplos}
\markboth{Exemplos}{Exemplos}

Tudo neste capítulo foi \textbf{executado}, não estimado. Os números são os
medidos, e cada um tem um script que o reproduz.

\section{O teste que você roda primeiro}

\subsection{Dois sabores no vácuo}

O caso com resposta analítica. Compare a saída do propagador contra
$\sin^2 2\theta \sin^2(kL)$ e exija concordância de máquina.

\begin{center}
\begin{tabular}{@{}>{\sffamily\small}p{6.6cm} >{\small}l@{}}
\toprule
\sffamily\bfseries\small Cenário & \sffamily\bfseries\small Erro medido \\
\midrule
2 sabores, vácuo, contra fórmula fechada & $2{,}1\times10^{-15}$ \\
3 sabores \emph{com} fase CP, em matéria, contra \cmd{expm} do SciPy & $7{,}2\times10^{-15}$ \\
contra \cmd{NuOscProbExact}, código público de método diferente & pior caso $2{,}2\times10^{-14}$ em 320 comparações \\
\bottomrule
\end{tabular}
\end{center}

O terceiro é o que vale mais. Os dois primeiros comparam o código contra a
matemática e contra si mesmo por outra rota; o terceiro compara contra uma
implementação independente, escrita por outras pessoas com outro método.

\begin{alerta}
\rotuloa{cuidado com a segunda implementação}
Os cinco scripts antigos que montavam a PMNS e o hamiltoniano em Python foram
\textbf{aposentados}. Eram uma segunda implementação da mesma física, e uma
segunda implementação só concorda com a primeira até o dia em que uma das duas
muda. O único cálculo em Python que sobrou é a fórmula analítica de referência,
que precisa ser independente \emph{de propósito}.
\end{alerta}

\subsection{A ressonância MSW}

Com o ângulo fixo, ligue o potencial de matéria e varra a densidade. A
amplitude de conversão sobe de $0{,}117$ para $0{,}999$ --- sem tocar em
$\theta$. É o efeito inteiro numa linha de saída: matéria não muda a mistura,
muda o problema.

\subsection{Um envelope de supernova}

Para $\rho \propto r^{-3}$, o GHOST entrega, num único \cmd{ghost run}:
\[
  P_{\rm inc}(\nue) = 0{,}0223635
  \qquad\text{contra}\qquad
  \sin^2\theta_{13} = 0{,}0223608,
\]
que é o valor adiabático esperado. A quinta casa.

\subsection{Custo}

Um espectro completo do DUNE, 200 energias $\times$ 2 sinais: \textbf{13\,ms}.
Isto importa mais do que parece --- um propagador que custa minutos não entra
num ajuste, e um que custa milissegundos entra.

\section{Validação contra a única conversão MSW que a natureza mediu}

Os testes acima comparam contra matemática e contra outro código. Nenhum
compara contra \emph{dado}. E existe exatamente uma conversão MSW medida: a do
Sol.

A comparação é feita canal a canal contra o Borexino, que mediu a
probabilidade de sobrevivência em várias energias --- do pp em
$0{,}267$~MeV, onde o código dá $0{,}54$, até o ${}^{8}$B, já no regime de
matéria. Não reproduzo os valores aqui porque eles mudam quando a tabela de
referência muda, e uma
apostila que cita número de validação desatualizado é pior que uma que não
cita.

\begin{nota}
\rotulo{por que este teste é diferente}
Os neutrinos solares de baixa energia saem no regime de vácuo, e os de alta
energia no regime de matéria adiabática. A transição entre os dois regimes ---
o ``sobe-desce'' da curva de sobrevivência --- é a assinatura do MSW, e é a
única medida direta dele. Um código que reproduz a fórmula fechada mas erra a
curva solar tem um problema de física, não de numérica.
\end{nota}

\section{A pergunta que o código foi construído para responder}

Uma supernova real não tem perfil suave. Há convecção, instabilidade do choque
de acreção, turbulência numa ampla faixa de escalas, e uma frente de choque
movendo-se para fora a dez mil quilômetros por segundo.

O GHOST estabeleceu \textbf{duas condições} que dizem quais dessas estruturas a
ressonância consegue sentir.

\subsection{Primeira condição: a escala}

Excite um modo de Fourier por vez sobre um perfil suave e meça o vazamento
$\Delta = P_{\rm inc}(\nue) - \sin^2\theta_{13}$, que se anula no limite
adiabático. A resposta é \emph{zero} para $\lambda > \Losc$ e cai como
$\lambda^{1{,}4\pm0{,}1}$ abaixo disso. Cinco ordens de grandeza em escala.

\subsection{Segunda condição: a amplitude}

Com a escala dentro da janela, o vazamento cresce como $\sigma^{2{,}02\pm0{,}03}$
sobre um fator 33 em amplitude, com o coeficiente reduzido $\Delta/\sigma^2$
variando apenas 9\%. Lei quadrática, saturando em um contraste de densidade de
$\sigma \simeq 8\%$.

\begin{figure}[!htb]
\centering
\begin{tikzpicture}
\begin{semilogxaxis}[
  width=12.4cm, height=5.2cm,
  xlabel={amplitude do contraste de densidade $\sigma$},
  ylabel={vazamento $\Delta$ (u.a.)},
  xmin=3e-4, xmax=3e-1, ymin=0, ymax=1.18,
  grid=major, grid style={linha, line width=0.4pt},
  tick label style={font=\sffamily\scriptsize, color=grafite},
  label style={font=\sffamily\small, color=tinta},
  axis line style={grafite},
]
% lei quadratica, saturando em 8%
\addplot[cyanescuro, line width=1.4pt, domain=3e-4:8e-2, samples=120]
  {(x/0.08)^2};
\addplot[cyanescuro, line width=1.4pt, domain=8e-2:3e-1, samples=40]{1};
\draw[goldescuro, densely dashed, line width=0.9pt]
  (axis cs:0.08,0) -- (axis cs:0.08,1.18);
\node[anchor=south east, font=\sffamily\scriptsize, text=goldescuro, rotate=90]
  at (axis cs:0.077,0.10) {saturação, $\sigma \simeq 8\%$};
\node[anchor=west, font=\sffamily\scriptsize, text=cyanescuro]
  at (axis cs:1.2e-3,0.55) {$\Delta \propto \sigma^{2{,}02\pm0{,}03}$};
% onde caem as estruturas reais
\draw[grafite, dotted, line width=0.8pt] (axis cs:7e-4,0) -- (axis cs:7e-4,1.18);
\node[anchor=west, font=\sffamily\scriptsize, text=grafite, rotate=90]
  at (axis cs:6.6e-4,0.10) {turbulência real};
\end{semilogxaxis}
\end{tikzpicture}
\caption{A segunda condição. A turbulência que os modelos de fato exibem está
$111\times$ abaixo da saturação; a frente de choque está $50\times$ acima. Não
é uma questão de grau --- as duas caem em regimes diferentes.}
\end{figure}

\subsection{O que as duas condições classificam}

\begin{center}
\begin{tabular}{@{}>{\sffamily\small}p{5.0cm} >{\small}c >{\small}c >{\small}c@{}}
\toprule
\sffamily\bfseries\small Estrutura & \sffamily\bfseries\small escala $<\Losc$?
  & \sffamily\bfseries\small amplitude $\gtrsim 8\%$? & \sffamily\bfseries\small importa? \\
\midrule
turbulência, $\sigma = 7\times10^{-4}$ & sim & não, $111\times$ abaixo & não \\
convecção, SASI, fluxos grandes & não, grandes demais & sim & não \\
\textbf{frente de choque} & \textbf{sim, é descontinuidade} & \textbf{sim, $50\times$ acima} & \textbf{sim} \\
\bottomrule
\end{tabular}
\end{center}

Uma estrutura precisa passar nas \emph{duas} peneiras. Só a frente de choque
passa.

\section{E aí vem o problema prático}

A frente qualifica porque é uma descontinuidade. A questão passa a ser: que
largura uma simulação hidrodinâmica de fato entrega, medida em comprimentos de
oscilação na ressonância?

\begin{figure}[!htb]
\centering
\begin{tikzpicture}
\begin{semilogxaxis}[
  width=12.4cm, height=5.0cm,
  xlabel={largura da frente, em unidades de $\Losc$},
  ylabel={$\PH$},
  xmin=0.02, xmax=20, ymin=0, ymax=1.05,
  grid=major, grid style={linha, line width=0.4pt},
  tick label style={font=\sffamily\scriptsize, color=grafite},
  label style={font=\sffamily\small, color=tinta},
  axis line style={grafite},
]
\addplot[cyanescuro, line width=1.4pt, domain=0.02:20, samples=200]
  {1/(1 + (x/0.55)^2.4)};
\draw[goldescuro, densely dashed, line width=0.9pt] (axis cs:0.1,0) -- (axis cs:0.1,1.05);
\draw[goldescuro, densely dashed, line width=0.9pt] (axis cs:3,0) -- (axis cs:3,1.05);
\node[anchor=south west, font=\sffamily\scriptsize, text=cyanescuro]
  at (axis cs:0.022,0.86) {satura: efeito total};
\node[anchor=south east, font=\sffamily\scriptsize, text=grafite]
  at (axis cs:19,0.06) {morre};
\node[anchor=south, font=\sffamily\scriptsize, text=goldescuro] at (axis cs:0.55,0.55)
  {a janela};
\end{semilogxaxis}
\end{tikzpicture}
\caption{Abaixo de $0{,}1\,\Losc$ o efeito satura; acima de $3\,\Losc$ ele
morre. Uma frente numérica de duas a três zonas num modelo 3D mede
$5{,}9\,\Losc$ --- do lado errado.}
\end{figure}

\begin{alerta}
\rotuloa{o achado central, e é uma correção que caducou}
A frente de choque \textbf{precisa ser afiada à mão} antes de a evolução de
sabor ser calculada. Isso era prática padrão entre 2002 e 2013. \emph{Nenhum
pipeline atual faz isso.}

\medskip
Omitir o afiamento remove cerca de metade de $\PH$: de $0{,}31$ para $0{,}81$ em
valor absoluto, conforme o caso. E há um agravante medido: a \textbf{média
angular destrói a frente por um fator 18}. O arquivo público distribui perfis
promediados; cada raio individual cruza abruptamente, a média não.
\end{alerta}

\section{O que isso custa ao DUNE}

Um modelo de fonte real através do SNOwGLoBES, \cmd{ar40kt} a 10\,kpc: a
contagem de corrente carregada de $\nue$ cai de $1241{,}8$ para $1187{,}0$
eventos. O viés fica entre $-3{,}9\%$ e $-4{,}4\%$ conforme a família de fonte.

Poucos por cento, com sinal que segue a fonte. Não é um erro dramático --- mas
hoje o pipeline atribui \emph{erro zero} a esse passo, porque aplica uma
constante, e uma constante não tem barra de erro.
